\begin{abstract}
Data-driven surrogates of finite element (FE) models can accelerate repeated structural assessment, but their member forces usually violate equilibrium and they give no indication of when they should not be trusted. This paper develops equilibrium-projected graph neural network surrogates for the first-order linear analysis of planar moment frames, with every analysis run in the open-source PyNite library and verified against an independent stiffness solver, OpenSeesPy and closed-form solutions. A graph network ensemble predicts all nodal displacements and member forces in a dimensionless, load-case-separated form, so that superposition, homogeneity in the loads, the $1/E$ scaling of displacements and the boundary conditions hold exactly. A projection in a variance-weighted metric maps the predicted forces to the nearest field in static equilibrium. Two indicators are used for two different decisions: at deployment, a domain check on geometry and load-pattern descriptors followed by a split-conformal bound on the ensemble spread decides which frames are referred back to the FE solver; during training, the equilibrium residual of the raw prediction decides which frames receive new FE analyses. On 7{,}500 PyNite analyses and five independent training runs, the median end-moment error was $2.37\pm0.03\%$ on in-distribution frames and $1.99\pm0.08\%$ on design-driven frames, four to nine times lower than five tabular baselines given the same features; a StructGNN-style network reached similar accuracy, and displacement-only networks with force recovery were much less accurate. On unseen geometries and load patterns the median moment errors rose to 7.5\% and 8.7\%, with 95th percentiles above 20\%. The projection reduced nodal out-of-balance forces from 4--16\% to below $10^{-12}$ and the axial-force errors by factors of 2--5, but changed the moment errors only modestly. With thresholds fixed on validation data, the referral rule accepted 98\% of in-distribution frames with 1.4\% of them above a 10\% moment error, and referred most unseen frames. Residual-guided acquisition reached the accuracy of 1{,}200 random analyses on unseen geometries with 600 analyses, while variance-guided acquisition was worse than random sampling.
\end{abstract}
